\documentclass{situationist-preprint}

\title{The Indispensability Argument}
\msid{MTH-01-b}
\msfamily{MTH-01}
\msdiscipline{Philosophy of Mathematics}
\mstype{Companion manuscript}
\mskeywords{indispensability; Quine; Putnam; ontological commitment; confirmational holism; mathematical explanation}

\begin{document}
\maketitle

\begin{abstract}
MTH-01 holds that numbers are abstract objects and that the indispensability argument establishes their existence. We set out the argument in a form that makes its two premises separable: a criterion of ontological commitment for accepted theories, and a claim of indispensability for the mathematics used in physics. We show that the argument needs only a weak, comparative form of confirmational holism, and that this form survives the objections of Sober and Maddy, which target a stronger holism than the argument requires. We also show that the explanatory version of the argument, due to Baker, does not depend on holism at all, and so gives the conclusion a second, independent support.
\end{abstract}

\section{The argument}\label{s:arg}

The argument goes back to \citet{Quine1948} and \citet{Putnam1971}; its standard modern statement is that of \citet{Colyvan2001}. We give it in three steps.

\begin{principle}[Commitment]\label{p:commit}
If we are justified in accepting a theory, and the theory cannot be reformulated without quantifying over entities of a kind $K$ while retaining its empirical and explanatory virtues, we are justified in believing that entities of kind $K$ exist.
\end{principle}

\begin{assumption}[Indispensability]\label{a:indisp}
Our best physical theories quantify over numbers, and cannot be reformulated without doing so while retaining their virtues.
\end{assumption}

\begin{proposition}\label{p:exist}
We are justified in believing that numbers exist, and that they are abstract objects.
\end{proposition}

\begin{proof}
By Assumption~\ref{a:indisp} and Principle~\ref{p:commit}, we are justified in believing that numbers exist. The theories in question locate their concrete objects in space and time and assign them causal powers; they assign numbers neither location nor causal role, and represent them as the values of magnitudes of every concrete system indifferently. The same justification that supports their existence therefore supports their being abstract.
\end{proof}

Assumption~\ref{a:indisp} is the premise MTH-16 defends: Field's programme of nominalizing physics fails, and with it the main attempt to dispense with numbers. Principle~\ref{p:commit} is the premise to be examined here. It is sometimes stated as a claim of \emph{confirmational holism}: the evidence for a theory confirms every part of it, the mathematics included. Principle~\ref{p:commit} requires less.

\begin{definition}\label{d:weak}
\emph{Weak holism} is the thesis that if a hypothesis is an ineliminable part of a theory that is confirmed relative to its available rivals, then belief in that hypothesis is justified to a degree not less than the theory's.
\end{definition}

Weak holism does not say that mathematics receives confirmation in the way a hypothesis about electrons does, by surviving tests that could have refuted it. It says that, where there is no rival that does without the hypothesis, the hypothesis cannot be held to a lower standard of belief than the theory of which it is an ineliminable part. That is enough for Principle~\ref{p:commit}: the condition that no reformulation dispenses with $K$ while retaining the theory's virtues is precisely the condition that there is no rival without the hypothesis that $K$s exist.

\section{Objections}\label{s:obj}

\begin{objection}[Sober: no contrastive confirmation]
Confirmation is contrastive. A hypothesis is confirmed by an observation only relative to a rival under which the observation was less probable. Every empirically serious rival in physics uses the same mathematics, so no observation favors the mathematics over anything, and the mathematics is not confirmed \citep{Sober1993}.
\end{objection}

\begin{reply}
The objection is correct about strong holism, and it does not touch Definition~\ref{d:weak}. Weak holism does not claim that observations raise the probability of the mathematics. It claims that the mathematics cannot rationally be believed to a lower degree than the physics it is part of. If every serious rival uses the same mathematics, the mathematics is not tested; but neither is there any theory, however poorly confirmed, that would allow us to believe the physics and disbelieve the mathematics. Sober's point explains why mathematical belief is not revised in response to failed predictions. It does not supply a reason to withhold it.
\end{reply}

\begin{objection}[Maddy: practice]
Scientists do not treat the mathematics of their theories as true. They use continuous fluid models while believing that fluids are made of molecules, and they treat infinite domains as idealizations. If the parts of a theory that scientists themselves treat as useful fictions do not carry commitment, the indispensability argument cannot read commitment off the mathematics \citep{Maddy1992}.
\end{objection}

\begin{reply}
The examples show that some \emph{physical} hypotheses are held as idealizations: that the fluid is continuous, that the plane is frictionless. What is treated as an idealization is a claim about the concrete system, and the mathematics is what describes the ideal system to which the actual one is compared. To say that an actual fluid approximates a continuous one is to quantify over the real-valued functions that describe the continuous one. The practice Maddy describes therefore relocates the commitment; it does not remove it. Principle~\ref{p:commit} applies to the comparison, and the comparison is indispensable to the theory's use.
\end{reply}

\begin{objection}[Melia: weaseling]
One can assert a theory and then take back its mathematical commitments, as one may say \enquote{everyone in the room is a spy, except the one I am talking to}. If the weaseled theory is as good as the original, Principle~\ref{p:commit} does not apply \citep{Melia2000}.
\end{objection}

\begin{reply}
The weaseled theory is not a reformulation that retains the virtues of the original; it is the original with a disclaimer whose content cannot be stated without the vocabulary it disclaims. The description of the concrete world it leaves standing is given only as \enquote{whatever the mathematical theory says about concrete things}, and that description has no content independent of the mathematics. A theory whose concrete content can be specified only by reference to the abstract objects it denies has not dispensed with them.
\end{reply}

\section{The explanatory version}\label{s:expl}

A second form of the argument does not rely on holism. \citet{Baker2005} argued that some physical phenomena have explanations in which a mathematical fact does the explanatory work: the prime-numbered life cycles of periodical cicadas are explained by the fact that prime periods minimize overlap with the cycles of predators of other periods, and the explanation turns on a number-theoretic property. If inference to the best explanation is warranted in science, and the best explanation of a physical phenomenon includes a mathematical fact essentially, we are warranted in believing that fact and in accepting what it quantifies over.

\begin{proposition}\label{p:expl}
Given inference to the best explanation, there is a justification for belief in numbers that does not depend on Definition~\ref{d:weak}.
\end{proposition}

\begin{proof}
The cicada explanation requires that the relevant periods have the property of being prime, and that a prime number has no divisors other than one and itself. The explanation is not available in a nominalist paraphrase that speaks only of cicadas and years, since what explains is a general fact about all periods, not a description of the actual ones. Inference to the best explanation therefore warrants belief in the fact, and the fact quantifies over numbers. No premise about the distribution of confirmation over the parts of a theory was used.
\end{proof}

The two versions are complementary. The holistic version shows that we cannot believe our physics and disbelieve its mathematics; the explanatory version shows that the mathematics does work in physics that only its truth could do. Each supports MTH-01's conclusion that numbers exist and are abstract, and each is independent of the a priori route by which, on EPI-12, arithmetic can also be known.

\section{Conclusion}

The indispensability argument needs only weak holism, and weak holism survives the objections from contrastive confirmation, from scientific practice and from weaseling. The explanatory version adds a route that needs no holism. Either way, the existence of numbers as abstract objects is established on the same footing as the physics in which they are used.

\begin{thebibliography}{9}\small
\bibitem[Baker(2005)]{Baker2005} Baker, A. (2005). Are there genuine mathematical explanations of physical phenomena? \emph{Mind}, 114(454), 223--238.
\bibitem[Colyvan(2001)]{Colyvan2001} Colyvan, M. (2001). \emph{The Indispensability of Mathematics}. Oxford University Press.
\bibitem[Maddy(1992)]{Maddy1992} Maddy, P. (1992). Indispensability and practice. \emph{The Journal of Philosophy}, 89(6), 275--289.
\bibitem[Melia(2000)]{Melia2000} Melia, J. (2000). Weaseling away the indispensability argument. \emph{Mind}, 109(435), 455--479.
\bibitem[Putnam(1971)]{Putnam1971} Putnam, H. (1971). \emph{Philosophy of Logic}. Harper and Row.
\bibitem[Quine(1948)]{Quine1948} Quine, W.~V. (1948). On what there is. \emph{The Review of Metaphysics}, 2(5), 21--38.
\bibitem[Sober(1993)]{Sober1993} Sober, E. (1993). Mathematics and indispensability. \emph{The Philosophical Review}, 102(1), 35--57.
\end{thebibliography}

\end{document}
